% English translation, made from our own transcription (original.tex), not from any existing
% translation (the Langton 1999 English version is in copyright and was NOT consulted). Every math
% expression, symbol, \tag{}, \origpage{N} marker, \rmfigure image path, and \ednote is reproduced
% unchanged; only the prose is translated — including the Latin connectives set inside display math
% via \text{...}, \text{et} → \text{and} and \text{seu} → \text{or} (HOUSESTYLE R21). The two
% flagged misprints (the p.~39 radical
% $\sqrt{1+y^{4}}$ and the §34 condition $\mathfrak{C}=0$) are kept exactly as printed, with their
% editorial \ednote{} notes carried over. Euler's own abbreviation "arc." and his Fraktur
% coefficients ($\mathfrak{A}..\mathfrak{E}$, distinct variables from roman A..E, per HOUSESTYLE
% R13) are preserved. "A sin." is Euler's notation for the arcsine.
\documentclass{article}
\usepackage{readmasters}
\begin{document}

\origpage{37}
\section*{On the integration of the differential equation}

\[ \frac{m\,dx}{\sqrt{1-x^{4}}} = \frac{n\,dy}{\sqrt{1-y^{4}}} \]

\emph{By the author L. EULER.}

\textbf{§.~1.} When first, on the occasion of the discoveries of the illustrious Count Fagnano, I
had contemplated this equation, I did indeed elicit an algebraic relation between the variables $x$
and $y$ of such a kind as would satisfy this equation; but that relation could not be taken for a
complete integral equation, for the reason that it did not embrace an arbitrary constant quantity,
of the kind that is always wont to be introduced into the calculation through integration. For
hence, as is well enough known, incomplete and particular integrals are wont to be distinguished,
of which the former exhaust the whole force of the differential equations, whereas the latter
satisfy them only in such a way that yet other expressions can equally satisfy them. But the
criterion of a complete integral equation consists in this, that it must involve a constant quantity
which does not appear in the differential equation.

\origpage{38}
\textbf{§.~2.} That these things may be perceived more clearly, it will suffice to have considered
the simplest differential equation $dx = dy$, which is certainly satisfied by this integral $x = y$;
yet in effect this integral extends less widely than the differential $dx = dy$, since the same is
equally satisfied by this integral $x = y + a$, which extends much more widely, taking for $a$ any
constant quantity whatever; and this last integral is judged at length to exhaust the whole force of
the differential equation $dx = dy$, whence too it is called the complete integral equation, because
there is present in it the constant quantity $a$, which does not occur in the differential equation.
But if in place of that indefinite constant $a$ determinate values be substituted, from the complete
integral there are obtained particular integrals, which for this very reason extend less widely than
the proposed differential equation.

\textbf{§.~3.} But often the particular algebraic integral of a differential equation can be
exhibited, while nevertheless the complete integral is transcendental; this indeed happens if the
transcendental part has been multiplied by that arbitrary constant, which therefore, when that
constant is set equal to nothing, vanishes from the calculation and leaves an algebraic particular
integral. Thus this equation $dy = dx + (y-x)\,dx$ is manifestly satisfied by the value $y = x$; yet
by this only a particular integral is contained, since the complete one is $y = x + ae^{x}$, $e$
denoting the number whose logarithm is $=1$. Unless therefore the arbitrary constant $a$ be set
vanishing, the integral will always be transcendental.

\origpage{39}
\textbf{§.~4.} Since therefore it can happen that a differential equation admits a particular
algebraic integral, even though the complete integral be transcendental, so too there is no lack of
grounds for doubting whether the complete integral of the proposed differential equation
$\dfrac{m\,dx}{\sqrt{1-x^{4}}} = \dfrac{n\,dy}{\sqrt{1-y^{4}}}$ involves transcendental quantities,
even though for it a particular algebraic integral has been able to be exhibited. For since the
complete integral is:\ednote{In the printed source the second radical of this display is set with a plus, $\sqrt{1+y^{4}}$; an apparent misprint for $\sqrt{1-y^{4}}$ --- the sign is negative in the differential equation itself and everywhere else. Reproduced here exactly as printed.}
\[ m\int \frac{dx}{\sqrt{1-x^{4}}} = n\int \frac{dy}{\sqrt{1+y^{4}}} + C \]
but these integrals can in no way be assigned by calling in aid the quadrature either of the circle
or of the hyperbola, it seems by no means probable that those formulas, so greatly transcendental in
general — in such a way that the constant $C$ remains indeterminate — can be brought back to an
algebraic relation between $x$ and $y$.

\textbf{§.~5.} It is indeed known that the complete integral of this differential equation
$\dfrac{m\,dx}{\sqrt{1-xx}} = \dfrac{n\,dy}{\sqrt{1-yy}}$ can always be exhibited algebraically,
provided the proportion of the coefficients $m$ and $n$ be rational; but because the integral of
each formula denotes an arc of a circle, so that the complete integral is
$m\,\mathrm{A}\,\mathrm{sin.}\,x = n\,\mathrm{A}\,\mathrm{sin.}\,y + C$, while the relation of the
sines that belong to arcs holding a rational proportion between themselves can be expressed
algebraically, it is no wonder that in these cases too the complete integral equation can be
exhibited algebraically. But since a comparison of this kind has no place in the transcendental
formulas $\displaystyle\int \frac{dx}{\sqrt{1-x^{4}}}$ and
$\displaystyle\int \frac{dy}{\sqrt{1-y^{4}}}$, or at least is not
\origpage{40}
established, thence the reduction of the integral to algebraic quantities cannot be sought.

\textbf{§.~6.} Nevertheless I have observed that, if a differential equation of this kind be proposed
\[ \frac{m\,dx}{\sqrt{1-x^{4}}} = \frac{n\,dy}{\sqrt{1-y^{4}}} \]
the complete integral too — that is, one which involves an arbitrary constant quantity — can always
be expressed algebraically, provided the ratio $m:n$ be rational: which indeed seems to me the more
worthy of note, in that I was led to this integral by no sure method, but rather elicited it by
trying, or by divining. Whence there is no doubt that a direct method, leading to this same
integral, would enlarge the bounds of analysis in no mean degree; and its investigation therefore
seems to be commended to Analysts with all zeal.

\textbf{§.~7.} But the complete integral of that differential equation, whatever be the rational
ratio of the coefficients $m$ and $n$, I have been able to derive from the complete integration of
this equation $\dfrac{dx}{\sqrt{1-x^{4}}} = \dfrac{dy}{\sqrt{1-y^{4}}}$: for, this being granted, I
shall point out a sure method of concluding from it also the complete integral of this much more
widely extending equation
$\dfrac{m\,dx}{\sqrt{1-x^{4}}} = \dfrac{n\,dy}{\sqrt{1-y^{4}}}$. Which method could also be employed in
general for finding the integrals of equations of this kind $mX\,dx = nY\,dy$, provided only the
complete integral of this $X\,dx = Y\,dy$ has been dug out, and $Y$ denote such a function of $y$ as
$X$ is of $x$.

\origpage{41}
\textbf{§.~8.} I shall begin, then, from this equation
\[ \frac{dx}{\sqrt{1-x^{4}}} = \frac{dy}{\sqrt{1-y^{4}}} \]
which at first glance is evidently satisfied by the equation $x = y$, which is therefore its
particular integral. But then the same equation is also satisfied by that algebraic value
$x = -\sqrt{\dfrac{1-yy}{1+yy}}$; for since
$dx = +\dfrac{2y\,dy}{(1+yy)\sqrt{(1-yy)(1+yy)}}$ and $\sqrt{1-x^{4}} = \dfrac{2y}{1+yy}$, there will
be $\dfrac{dx}{\sqrt{1-x^{4}}} = \dfrac{dy}{\sqrt{1-y^{4}}}$. Hence this value too, that is, the
equation $xxyy + xx + yy - 1 = 0$, is a particular integral of the proposed differential equation.
Whence the complete integral, which involves an arbitrary constant, must necessarily be so
constituted that, on assigning to this constant a certain definite value, there results $x = y$; but
if to the same constant another value be assigned, there results
$x = -\sqrt{\dfrac{1-yy}{1+yy}}$, that is, $xxyy + xx + yy - 1 = 0$.

\subsection*{Theorem.}
\textbf{§.~9.} I assert, then, that of this differential equation
\[ \frac{dx}{\sqrt{1-x^{4}}} = \frac{dy}{\sqrt{1-y^{4}}} \]
the complete integral equation is:
\[ xx + yy + ccxxyy = cc + 2xy\sqrt{1-c^{4}} \]

\subsection*{Demonstration.}
For, this equation being posited, its differential will be:
\[ x\,dx + y\,dy + ccxy(x\,dy + y\,dx) = (x\,dy + y\,dx)\sqrt{1-c^{4}} \]
\origpage{42}
whence there arises
\[ dx\bigl(x + ccxyy - y\sqrt{1-c^{4}}\bigr) + dy\bigl(y + ccxxy - x\sqrt{1-c^{4}}\bigr) = 0 \]
But from the same equation, solved, there is gathered:
\[ y = \frac{x\sqrt{1-c^{4}} + c\sqrt{1-x^{4}}}{1+ccxx} \quad\text{and}\quad
   x = \frac{y\sqrt{1-c^{4}} - c\sqrt{1-y^{4}}}{1+ccyy} \]
For if there the radical $\sqrt{1-x^{4}}$ is given the sign $+$, here the radical $\sqrt{1-y^{4}}$
must be given the sign $-$, so that, on setting $x = 0$, the same value $y = c$ results on both
sides. There will be, therefore,
\[ x + ccxyy - y\sqrt{1-c^{4}} = -c\sqrt{1-y^{4}} \]
\[ y + ccxxy - x\sqrt{1-c^{4}} = c\sqrt{1-x^{4}} \]
and, these values being substituted in the differential equation, there results
\[ -c\,dx\sqrt{1-y^{4}} + c\,dy\sqrt{1-x^{4}} = 0, \]
that is
\[ \frac{dx}{\sqrt{1-x^{4}}} = \frac{dy}{\sqrt{1-y^{4}}}. \]
The integral of this differential equation is therefore:
\[ xx + yy + ccxxyy = cc + 2xy\sqrt{1-c^{4}} \]
and because it contains the constant $c$, depending on our choice, it will at the same time be the
complete integral. Q.~E.~D.

\textbf{§.~10.} If therefore this equation
$\dfrac{dx}{\sqrt{1-x^{4}}} = \dfrac{dy}{\sqrt{1-y^{4}}}$ be had, the complete integral value of $x$
is:
\[ x = \frac{y\sqrt{1-c^{4}} + c\sqrt{1-y^{4}}}{1+ccyy} \]
whence, if the arbitrary constant $c$ vanish, there results $x = y$; but if $c = 1$ be set, we have
$x = \pm\dfrac{\sqrt{1-y^{4}}}{1+yy} = \sqrt{\dfrac{1-yy}{1+yy}}$, which are both those particular
values already exhibited above.

\origpage{43}
Hence are elicited other particular values, simpler than the rest, but which fall away into
imaginaries. Thus, on setting $c = \infty$, there results $x = \dfrac{\sqrt{-1}}{y}$; and on setting
$cc = -1$, there results
$x = \sqrt{\dfrac{yy+1}{yy-1}}$, which likewise satisfy the proposed equation.

\textbf{§.~11.} But that the rationale of this integral may be perceived more clearly, let a curve
$AM$ be conceived of such a character that, the abscissa $AP = u$ being set, the arc corresponding to
it is $AM = \displaystyle\int \frac{du}{\sqrt{1-u^{4}}}$.
\rmfigure{figures/fig-1.png}{Fig.~1.}{Fig. 1 (Tab. I): the curve AM with abscissa AP; the arc from A measures the elliptic integral treated in the paper. Cropped from the original engraved plate, Novi Commentarii vol. VI (1761).}
Then, the same curve being described anew, let the abscissa
$ap = x$ be taken; the arc will be $am = \displaystyle\int \frac{dx}{\sqrt{1-x^{4}}}$. Taking, then,
\[ x = \frac{u\sqrt{1-c^{4}} + c\sqrt{1-u^{4}}}{1+ccuu} \]
there will result $\dfrac{dx}{\sqrt{1-x^{4}}} = \dfrac{du}{\sqrt{1-u^{4}}}$; and therefore
arc.~$am = $ arc.~$AM + $
Const. But for the determination of this constant, on setting $u = 0$ — in which case the arc $AM$
vanishes — there results $x = c$. Wherefore, if the abscissa $ab = c$ be taken, to which the arc $ad$
corresponds, the arc $dm = $ the arc $AM$.

\textbf{§.~12.} By the help, therefore, of this complete integration of the equation
$\dfrac{dx}{\sqrt{1-x^{4}}} = \dfrac{du}{\sqrt{1-u^{4}}}$, in the proposed curve there can be cut off,
equal to any arc $AM$ whatever which corresponds to the abscissa $AP = u$, an arc $dm$ that begins
from a given point $d$. For, the abscissa $ab = c$ corresponding to the given point $d$ being set, if
the abscissa $ap = x = \dfrac{c\sqrt{1-u^{4}} + u\sqrt{1-c^{4}}}{1+ccuu}$ be taken, the arc $dm$ will
be equal to the arc $AM$.
\rmfigure{figures/fig-2.png}{Fig.~2.}{Fig. 2 (Tab. I): a second copy of the curve, with the abscissae ap and a-pi (points a, pi, b, p, d, m) by which an arc dm or d-mu equal to a given arc AM is cut off from the fixed point d. Cropped from the original engraved plate, Novi Commentarii vol. VI (1761).}
But in like manner, since $\sqrt{1-c^{4}}$ may be set negative, if the abscissa
\[ a\pi = \frac{c\sqrt{1-u^{4}} - u\sqrt{1-c^{4}}}{1+ccuu} \]
\origpage{44}
be taken, the arc $d\mu$ will likewise be equal to the arc $AM$: and thus in this curve, from any
given point $d$, arcs $dm$ and $d\mu$ that are equal to the arc $AM$ can be cut off on both sides.

\textbf{§.~13.} Hence, then, it is clear that if the arc $ad$ be taken equal to the arc $AM$, that is
$c = u$, the arc $am$ will be double the arc $AM$. Hence if $ap = x = \dfrac{2u\sqrt{1-u^{4}}}{1+u^{4}}$
be set, there will come out the arc $am = 2$~arc.~$AM$. In like manner, if the arc $ad = 2AM$ be
taken, that is $c = \dfrac{2u\sqrt{1-u^{4}}}{1+u^{4}}$, and
$x = \dfrac{c\sqrt{1-u^{4}} + u\sqrt{1-c^{4}}}{1+ccuu}$ be set, there will be obtained the arc
$am = 3$~arc.~$AM$. And if that value of $x$ be again substituted for $c$, so that $ad = 3AM$, and
again $x = \dfrac{c\sqrt{1-u^{4}} + u\sqrt{1-c^{4}}}{1+ccuu}$ be set, there will be born an arc $am$
quadruple the arc $AM$; and so on successively, any multiples whatever of the arc $AM$ will be able
to be assigned geometrically.

\textbf{§.~14.} Let the arc $ad = n \cdot AM$ and $ab = z$; so that
$\displaystyle\int \frac{dz}{\sqrt{1-z^{4}}} = n\int \frac{du}{\sqrt{1-u^{4}}}$; and from these it is clear that,
if $x = \dfrac{z\sqrt{1-u^{4}} + u\sqrt{1-z^{4}}}{1+uuzz}$ be taken, there will be
$\displaystyle\int \frac{dx}{\sqrt{1-x^{4}}} = (n+1)\int \frac{du}{\sqrt{1-u^{4}}}$; but if
$x = \dfrac{z\sqrt{1-u^{4}} - u\sqrt{1-z^{4}}}{1+uuzz}$ be set, then it will come about that
$\displaystyle\int \frac{dx}{\sqrt{1-x^{4}}} = (n-1)\int \frac{du}{\sqrt{1-u^{4}}}$. If therefore this equation
$\dfrac{dz}{\sqrt{1-z^{4}}} = \dfrac{n\,du}{\sqrt{1-u^{4}}}$ has been integrated, and the due value for
$z$ dug out therefrom, this equation too
$\dfrac{dx}{\sqrt{1-x^{4}}} = \dfrac{(n\pm1)\,du}{\sqrt{1-u^{4}}}$ will be able to be integrated, its
integral being of course
$x = \dfrac{z\sqrt{1-u^{4}} \pm u\sqrt{1-z^{4}}}{1+uuzz}$. And if for $z$ its complete value be
assumed — that is, one involving an arbitrary constant — for $x$ too its complete value will come
out.

\origpage{45}
\textbf{§.~15.} Hence, then, it is evident how the complete integral equation ought to be found that
agrees with this differential equation
$\dfrac{dx}{\sqrt{1-x^{4}}} = \dfrac{n\,du}{\sqrt{1-u^{4}}}$, as often as $n$ is a whole number. But in
like manner $y$ can be assigned, so that
$\dfrac{dy}{\sqrt{1-y^{4}}} = \dfrac{m\,du}{\sqrt{1-u^{4}}}$; whence if, by eliminating $u$, the
equation between $x$ and $y$ be sought, it will be the integral of this equation
$\dfrac{m\,dx}{\sqrt{1-x^{4}}} = \dfrac{n\,dy}{\sqrt{1-y^{4}}}$, whatever rational numbers be
substituted for $m$ and $n$: and that this integral may come out complete, it suffices to have
determined, for one only of the variables $x$ and $y$, the complete value through $u$, since hereby a
new arbitrary constant is now introduced into the calculation.

\textbf{§.~16.} The method that I have here used in the demonstration of the Theorem, although it is
not drawn from the nature of the thing, but led indirectly to what was proposed, nevertheless
extends much more widely: for in like manner it is gathered that the complete integral of this
differential equation
\[ \frac{dx}{\sqrt{1+mxx+nx^{4}}} = \frac{dy}{\sqrt{1+myy+ny^{4}}} \]
is:
\[ 0 = cc - xx - yy + nccxxyy + 2xy\sqrt{1+mcc+nc^{4}} \]
Whence, by applying the same reasoning as before, the complete integral of this equation too
\[ \frac{\mu\,dx}{\sqrt{1+mxx+nx^{4}}} = \frac{\nu\,dy}{\sqrt{1+myy+ny^{4}}} \]
will be obtained, provided that by the letters $\mu$ and $\nu$ whole numbers be designated.

\textbf{§.~17.} But the investigation of this integration proceeds thus: Let there first be feigned,
at will, a relation between the variables $x$ and $y$ contained in this equation:
\origpage{46}
\[ \alpha xx + \alpha yy = 2\beta xy + \gamma xxyy + \delta \tag{1} \]
which, differentiated, gives:
\[ \alpha x\,dx + \alpha y\,dy = \beta x\,dy + \beta y\,dx + \gamma xyy\,dx + \gamma xxy\,dy \]
whence there is made up
\[ dx(\alpha x - \beta y - \gamma xyy) + dy(\alpha y - \beta x - \gamma xxy) = 0 \tag{2} \]
Then from equation (1) let the values of each variable be elicited:
\[ x = \frac{\beta y + \sqrt{\alpha\delta + (\beta\beta - \alpha\alpha - \gamma\delta)yy + \alpha\gamma y^{4}}}{\alpha - \gamma yy} \]
\[ y = \frac{\beta x - \sqrt{\alpha\delta + (\beta\beta - \alpha\alpha - \gamma\delta)xx + \alpha\gamma x^{4}}}{\alpha - \gamma xx} \]
And hence we obtain:
\[ \alpha x - \beta y - \gamma xyy = \sqrt{\alpha\delta + (\beta\beta - \alpha\alpha - \gamma\delta)yy + \alpha\gamma y^{4}} \tag{3} \]
\[ \alpha y - \beta x - \gamma xxy = -\sqrt{\alpha\delta + (\beta\beta - \alpha\alpha - \gamma\delta)xx + \alpha\gamma x^{4}} \tag{4} \]
which values, substituted in equation (2), will afford
\[ \frac{dx}{\sqrt{\alpha\delta + (\beta\beta - \alpha\alpha - \gamma\delta)xx + \alpha\gamma x^{4}}}
   = \frac{dy}{\sqrt{\alpha\delta + (\beta\beta - \alpha\alpha - \gamma\delta)yy + \alpha\gamma y^{4}}} \tag{5} \]
the integral of which equation is therefore equation (1).

\textbf{§.~18.} That we may render those forms simpler, let us put $\alpha\delta = A$;
$\beta\beta - \alpha\alpha - \gamma\delta = C$; $\alpha\gamma = E$; and there will be $\delta = \dfrac{A}{\alpha}$;
$\gamma = \dfrac{E}{\alpha}$ and $\beta = \sqrt{C + \alpha\alpha + \dfrac{AE}{\alpha\alpha}}$. Wherefore of this
differential equation
\[ \frac{dx}{\sqrt{A + Cxx + Ex^{4}}} = \frac{dy}{\sqrt{A + Cyy + Ey^{4}}} \tag{6} \]
the integral equation is this:
\[ \alpha(xx + yy) = \frac{A}{\alpha} + \frac{E}{\alpha}xxyy + 2xy\sqrt{C + \alpha\alpha + \frac{AE}{\alpha\alpha}} \tag{7} \]
which is at the same time the complete integral:

\textbf{§.~19.} Or let us put $A = f\alpha\alpha$: $C = g\alpha\alpha$ and $E = b\alpha\alpha$, that we
may have this differential equation
\[ \frac{dx}{\sqrt{f + gxx + bx^{4}}} = \frac{dy}{\sqrt{f + gyy + by^{4}}} \]
\origpage{47}
of which therefore the complete integral equation will be:
\[ xx + yy = f + bxxyy + 2xy\sqrt{1 + g + fb} \]
which, although it does not seem to involve a new constant, is nevertheless complete, since in the
differential only the ratio of the quantities $f$, $g$, and $b$ is regarded, so that in place of
$f$, $g$, and $b$ one may write $fcc$,
$gcc$ and $bcc$, whence the integral equation comes out manifestly complete:
\[ xx + yy = fcc + bccxxyy + 2xy\sqrt{1 + gcc + fbc^{4}} \]
or $f(xx + yy) = fee + beexxyy + 2xy\sqrt{f(f + gee + be^{4})}$, setting $cc = \dfrac{ee}{f}$.

\textbf{§.~20.} But if, then, this differential equation be proposed
\[ \frac{dx}{\sqrt{f + gxx + bx^{4}}} = \frac{dy}{\sqrt{f + gyy + by^{4}}} \]
the value of $y$ will be able to be expressed through an algebraic function of $x$, so as to be:
\[ y = \frac{x\sqrt{1 + gcc + fbc^{4}} \pm c\sqrt{f + gxx + bx^{4}}}{1 - bccxx} \]
or $y = \dfrac{x\sqrt{f(f + gee + be^{4})} \pm e\sqrt{f(f + gxx + bx^{4})}}{f - beexx}$.

But if, then, $g = 0$, so that this differential equation be had
\[ \frac{dx}{\sqrt{f + bx^{4}}} = \frac{dy}{\sqrt{f + by^{4}}} \]
the complete integral value of $y$ will be
\[ y = \frac{x\sqrt{f(f + be^{4})} \pm e\sqrt{f(f + bx^{4})}}{f - beexx} \]
whence, by determining the constant $e$ at pleasure, innumerable particular values for $y$ can be
deduced.

\origpage{48}
\textbf{§.~21.} But by the benefit of the method that I used above, of this equation too
\[ \frac{m\,dx}{\sqrt{f + gxx + bx^{4}}} = \frac{n\,dy}{\sqrt{f + gyy + by^{4}}} \]
if only $m$ and $n$ be rational numbers, the complete integral — and that indeed algebraically —
will be able to be exhibited.

\textbf{§.~22.} Just as in the equation assumed above the variables $x$ and $y$ were made
interchangeable with one another, so that both formulas turned out similar to each other, so, this
limitation being omitted, we shall arrive at the comparison of unlike differential formulas. Let us
put, then:
\[ \alpha xx + \beta yy = 2\gamma xy + \delta xxyy + \varepsilon \tag{1} \]
whence there arises
\[ x = \frac{\gamma y + \sqrt{\alpha\varepsilon + (\gamma\gamma - \delta\varepsilon - \alpha\beta)yy + \beta\delta y^{4}}}{\alpha - \delta yy} \]
and $y = \dfrac{\gamma x - \sqrt{\beta\varepsilon + (\gamma\gamma - \delta\varepsilon - \alpha\beta)xx + \alpha\delta x^{4}}}{\beta - \delta xx}$
and hence
\[ \alpha x - \gamma y - \delta xyy = \sqrt{\alpha\varepsilon + (\gamma\gamma - \delta\varepsilon - \alpha\beta)yy + \beta\delta y^{4}} \tag{2} \]
\[ \beta y - \gamma x - \delta xxy = -\sqrt{\beta\varepsilon + (\gamma\gamma - \delta\varepsilon - \alpha\beta)xx + \alpha\delta x^{4}} \tag{3} \]
but equation (1), differentiated, gives:
\[ dx(\alpha x - \gamma y - \delta xyy) + dy(\beta y - \gamma x - \delta xxy) = 0 \]
whence there is made up this differential equation:
\[ \frac{dx}{\sqrt{\beta\varepsilon + (\gamma\gamma - \delta\varepsilon - \alpha\beta)xx + \alpha\delta x^{4}}}
   = \frac{dy}{\sqrt{\alpha\varepsilon + (\gamma\gamma - \delta\varepsilon - \alpha\beta)yy + \beta\delta y^{4}}} \]
the integral of which is therefore the assumed equation.

\textbf{§.~23.} But this disparity is easily removed by putting, in place of $y$,
$z\sqrt{\dfrac{\alpha}{\beta}}$,
the reason for which could have been at once manifest from the assumed equation. But another way
lies open for arriving at unlike formulas, of which it will suffice to have here handed down an
example. Let the equation be assumed:
$x^{4} + 2axxyy + 2bxx$
\origpage{49}
${}= c$, whose differential is $dx(x^{3} + axyy + bx) + axxy\,dy = 0$, that is
\[ \frac{dx}{xy} = \frac{-a\,dy}{xx + ayy + b} \]
Now from the assumed equation let first $xy$ be determined through $x$, and thus there will be
$xy = \sqrt{\dfrac{c - 2bxx - x^{4}}{2a}}$; but then $xx + ayy + b$ through $y$, and, because
$(xx + ayy + b)^{2} = c + (ayy + b)^{2}$, there will be
\[ xx + ayy + b = \sqrt{c + (ayy + b)^{2}} \]
Wherefore there will be had this differential equation
\[ \frac{dx\sqrt{2a}}{\sqrt{c - 2bxx - x^{4}}} = \frac{a\,dy}{\sqrt{c + bb + 2abyy + aay^{4}}} \]
the integral of which is therefore the assumed one, that is $y = \dfrac{\sqrt{c - 2bxx - x^{4}}}{x\sqrt{2a}}$.

\textbf{§.~24.} Although this integral is not complete, yet from what precedes it will easily be
rendered complete. For let there be put:
\[ \frac{a\,dy}{\sqrt{c + bb + 2abyy + aay^{4}}} = \frac{a\,dz}{\sqrt{c + bb + 2abzz + aaz^{4}}} \]
because $f = c + bb$; $g = 2ab$; $b = aa$, there will be
\[ y = \frac{z\sqrt{(c+bb)(c + bb + 2abee + aae^{4})} \pm e\sqrt{(c+bb)(c + bb + 2abzz + aaz^{4})}}{c + bb - aaeezz} \]
let this value, then, be set equal to $\dfrac{\sqrt{c - 2bxx - x^{4}}}{x\sqrt{2a}}$, and the equation
between $x$ and $z$ resulting hence will be the complete integral of this differential equation
\[ \frac{dx\sqrt{2a}}{\sqrt{c - 2bxx - x^{4}}} = \frac{a\,dz}{\sqrt{c + bb + 2abzz + aaz^{4}}} \]
Nay more, from what has been adduced it is clear, if these two members be moreover multiplied by any
rational numbers whatever, in what way the complete integral ought to be found.

\origpage{50}
\textbf{§.~25.} But, dismissing the disparity of the members, let us conceive the formation of equal
members more generally; let there be put, then:
\[ 0 = \alpha + 2\beta(x+y) + \gamma(xx+yy) + 2\delta xy + 2\varepsilon xy(x+y) + \zeta xxyy \tag{1} \]
whence, by differentiating, there is obtained:
\[ dx(\beta + \gamma x + \delta y + 2\varepsilon xy + \varepsilon yy + \zeta xyy)
   + dy(\beta + \gamma y + \delta x + 2\varepsilon xy + \varepsilon xx + \zeta xxy) = 0 \]
and therefore
\[ \frac{dy}{\beta + \gamma x + \delta y + 2\varepsilon xy + \varepsilon yy + \zeta xyy}
   = \frac{-dx}{\beta + \gamma y + \delta x + 2\varepsilon xy + \varepsilon xx + \zeta xxy} \tag{2} \]
But from the solution of the assumed equation there is elicited.
\[ y = \frac{-\beta - \delta x - \varepsilon xx + \sqrt{\beta\beta - \alpha\gamma + 2(\beta\delta - \alpha\varepsilon - \beta\gamma)x + (\delta\delta - \gamma\gamma - \alpha\zeta - 2\beta\varepsilon)xx + 2(\delta\varepsilon - \beta\zeta - \gamma\varepsilon)x^{3} + (\varepsilon\varepsilon - \gamma\zeta)x^{4}}}{\gamma + 2\varepsilon x + \zeta xx} \]
Let there be put, for brevity's sake,
\[ \beta\beta - \alpha\gamma = A; \quad \beta\delta - \alpha\varepsilon - \beta\gamma = B; \quad
   \delta\delta - \gamma\gamma - \alpha\zeta - 2\beta\varepsilon = C; \quad
   \varepsilon\varepsilon - \gamma\zeta = E; \quad \delta\varepsilon - \beta\zeta - \gamma\varepsilon = D; \]
and there will be
\[ \beta + \delta x + \varepsilon xx + \gamma y + 2\varepsilon xy + \zeta xxy = \pm\sqrt{A + 2Bx + Cxx + 2Dx^{3} + Ex^{4}} \]
\[ \beta + \delta y + \varepsilon yy + \gamma x + 2\varepsilon xy + \zeta xyy = \mp\sqrt{A + 2By + Cyy + 2Dy^{3} + Ey^{4}} \]

\textbf{§.~26.} Hence, therefore, we conclude that of this differential equation:
\[ \frac{dx}{\sqrt{A + 2Bx + Cxx + 2Dx^{3} + Ex^{4}}} = \frac{dy}{\sqrt{A + 2By + Cyy + 2Dy^{3} + Ey^{4}}} \]
the integral equation, and that a complete one, is
\[ 0 = \alpha + 2\beta(x+y) + \gamma(xx+yy) + 2\delta xy + 2\varepsilon xy(x+y) + \zeta xxyy \]
\origpage{51}
with, of course, the above determination of these coefficients being applied. But first let $\beta$
or $\varepsilon$ be defined from this equation
\[ \frac{BB(\varepsilon\varepsilon - E) - DD(\beta\beta - A)}{A\varepsilon\varepsilon - E\beta\beta}
   + \frac{2AD\varepsilon - 2BE\beta}{B\varepsilon - D\beta} = C \]
and then there will be:
\[ \gamma = \frac{A\varepsilon\varepsilon - E\beta\beta}{B\varepsilon - D\beta}; \quad
   \alpha = \frac{\beta\beta - A}{\gamma}; \quad \zeta = \frac{\varepsilon\varepsilon - E}{\gamma} \quad\text{and} \]
\[ \delta = \frac{B\beta(\varepsilon\varepsilon - E) - D\varepsilon(\beta\beta - A)}{A\varepsilon\varepsilon - E\beta\beta} + \gamma
   \quad\text{or}\quad \delta = \gamma + \frac{B + \alpha\varepsilon}{\beta} \]

\textbf{§.~27.} Hence, then, it is evident that this differential equation too:
\[ \frac{dx}{\sqrt{A + 2Dx^{3}}} = \frac{dy}{\sqrt{A + 2Dy^{3}}} \]
can be integrated: for, because $B = 0$, $C = 0$ and $E = 0$, there will be
\[ -\frac{DD(\beta\beta - A)}{A\varepsilon\varepsilon} - \frac{2A\varepsilon}{\beta} = 0 \quad\text{or}\quad
   \varepsilon = \sqrt[3]{\frac{DD}{2AA}\beta(A - \beta\beta)} \]
but hence the values come out too complicated. The business will be dispatched more easily by
resolving the values of the vanishing letters $B$, $C$ and $E$; for $E = 0$ gives: $\zeta = \dfrac{\varepsilon\varepsilon}{\gamma}$;
then $B = 0$ gives: $\delta = \gamma + \dfrac{\alpha\varepsilon}{\beta}$; and $C = 0$ gives
$\delta\delta - \gamma\gamma = \alpha\zeta + 2\beta\varepsilon = \dfrac{\alpha\varepsilon\varepsilon}{\gamma} + 2\beta\varepsilon
= \dfrac{\alpha^{2}\varepsilon\varepsilon}{\beta\beta} + \dfrac{2\alpha\gamma\varepsilon}{\beta}$, whose factors are
$\beta\beta = \alpha\gamma$ and $\alpha\varepsilon\varepsilon + 2\beta\gamma\varepsilon = 0$. But if
$\beta\beta = \alpha\gamma$ were the case, $A = 0$ would follow; while if $\varepsilon = 0$ were the case, both $\zeta = 0$ and
$D = 0$ would follow, contrary to the aim. It must therefore come about that $\alpha\varepsilon = -2\beta\gamma$; whence there will be
$\alpha = -\dfrac{2\beta\gamma}{\varepsilon}$; $\delta = -\gamma$; and $\zeta = \dfrac{\varepsilon\varepsilon}{\gamma}$.
Finally it must be that $\beta\beta + \dfrac{2\beta\gamma\gamma}{\varepsilon} = A$ and
$-2\gamma\varepsilon - \dfrac{\beta\varepsilon\varepsilon}{\gamma} = D$. Thence there arises
$\varepsilon = \dfrac{2\beta\gamma\gamma}{A - \beta\beta}$; and because
$\dfrac{\gamma D}{\varepsilon} = -(2\gamma\gamma + \beta\varepsilon)$ and $2\gamma\gamma + \beta\varepsilon = \dfrac{A\varepsilon}{\beta}$,
there will be $\dfrac{\gamma D}{\varepsilon} = -\dfrac{A\varepsilon}{\beta}$; and therefore $\varepsilon\varepsilon = -\dfrac{\beta\gamma D}{A}$. Therefore
\[ \frac{4\beta\gamma^{3}}{(A - \beta\beta)^{2}} + \frac{D}{A} = 0. \]

\origpage{52}
\textbf{§.~28.} But since only the ratio of the letters $A$ and $D$ comes into the reckoning, the
last equation serves for finding the absolute value of $A$, which however it is not needful to know.
The letters $\gamma$ and $\beta$ will therefore remain indeterminate. Let there be put, then, $\gamma = -Ac$ and $\beta = Dc$; there will be
$\varepsilon\varepsilon = DDcc$, that is $\varepsilon = Dc$, and hence $\delta = Ac$;
$\zeta = -\dfrac{DDc}{A}$; and $\alpha = 2Ac$. Wherefore of this differential equation:
\[ \frac{dx}{\sqrt{A + 2Dx^{3}}} = \frac{dy}{\sqrt{A + 2Dy^{3}}} \]
the integral is.
\[ 0 = 2A + 2D(x+y) - A(xx+yy) + 2Axy + 2Dxy(x+y) - \frac{DD}{A}xxyy \]
But this integral is not complete; it will, however, be rendered such by putting $\gamma = -A$ and $\beta = Dcc$,
whence there arises $\varepsilon\varepsilon = DDcc$ and $\varepsilon = Dc$; further there will be $\delta = A$;
$\zeta = -\dfrac{DDcc}{A}$; $\alpha = 2Ac$; so that the complete integral is:
\[ 0 = 2Ac + 2Dcc(x+y) - A(xx+yy) + 2Axy + 2Dcxy(x+y) - \frac{DDcc}{A}xxyy \]
where $c$ is a constant depending on choice, whence there arises:
\[ y = \frac{Dcc + Ax + Dcxx + \sqrt{c\bigl(2A + \frac{DD}{A}c^{3}\bigr)(A + 2Dx^{3})}}{A - 2Dcx + \frac{DDcc}{A}xx} \]

\textbf{§.~29.} This case deserves to be noted, in which $A = 1$ and $D = \frac{1}{2}$, so that this
differential equation be had
\[ \frac{dx}{\sqrt{1 + x^{3}}} = \frac{dy}{\sqrt{1 + y^{3}}} \]
where, to remove the fractions, let $2c$ be written in place of $c$, and the complete integral will
be:
\[ 0 = 4c + 4cc(x+y) - xx - yy + 2xy + 2cxy(x+y) - ccxxyy \]
that is $y = \dfrac{2cc + x + cxx + 2\sqrt{c(1+c^{3})(1+x^{3})}}{1 - 2cx + ccxx}$

\origpage{53}
The particular integrals, therefore, will be

I.~if $c = 0$; $y = x$;

II.~if $c = \infty$; $y = \dfrac{2 \pm 2\sqrt{1+x^{3}}}{xx}$;

III.~if $c = -1$; $y = \dfrac{2 + x - xx}{1 + 2x + xx} = \dfrac{2-x}{1+x}$.

\textbf{§.~30.} From the same principle, if in §.~29. in place of the letters $A$, $B$, $C$, $D$, $E$,
the same be multiplied by some quantity $p$, none the less the differential equation will be
\[ \frac{dx}{\sqrt{A + 2Bx + Cxx + 2Dx^{3} + Ex^{4}}} = \frac{dy}{\sqrt{A + 2By + Cyy + 2Dy^{3} + Ey^{4}}} \]
and there will be found
\[ p = \frac{BB\varepsilon\varepsilon - DD\beta\beta}{BBE - ADD}
   + 2\frac{(AD\varepsilon - BE\beta)(A\varepsilon\varepsilon - E\beta\beta)}{(B\varepsilon - D\beta)(BBE - ADD)}
   - \frac{C(A\varepsilon\varepsilon - E\beta\beta)}{BBE - ADD} \]
then there will be $\gamma = \dfrac{A\varepsilon\varepsilon - E\beta\beta}{B\varepsilon - D\beta}$;
$\alpha = \dfrac{\beta\beta - Ap}{\gamma}$; $\zeta = \dfrac{\varepsilon\varepsilon - Ep}{\gamma}$ and
$\delta = \gamma + \dfrac{\alpha\varepsilon + Bp}{\beta}$: so that the letters $\beta$ and $\varepsilon$ remain
indeterminate, and there will therefore be made the complete integral equation:
\[ 0 = \alpha + 2\beta(x+y) + \gamma(xx+yy) + 2\delta xy + 2\varepsilon xy(x+y) + \zeta xxyy \]
whence there arises:
\[ y = \frac{-\beta - \delta x - \varepsilon xx + \sqrt{p(A + 2Bx + Cxx + 2Dx^{3} + Ex^{4})}}{\gamma + 2\varepsilon x + \zeta xx} \]

\textbf{§.~31.} Finally it is to be noted that not only this differential equation, whose complete
integral I have just exhibited, but also this much more widely extending one
\[ \frac{m\,dx}{\sqrt{A + 2Bx + Cxx + 2Dx^{3} + Ex^{4}}} = \frac{n\,dy}{\sqrt{A + 2By + Cyy + 2Dy^{3} + Ey^{4}}} \]
can always be integrated algebraically, and indeed completely, provided the ratio of the
coefficients $m$ and $n$ be rational: for this integration is set up in a manner similar to that
which I used above for integrating the equation that had here been especially proposed to me. But
the method, of which
\origpage{54}
I have here brought forward specimens, seems to me so constituted that, by cultivating its character
more diligently, it can be rendered fit for signal uses, whence advantages by no means to be
despised will redound to Analysis.

\textbf{§.~32.} But here I observe that, by extending the formula assumed in §.~28 more widely,
differentials of such a kind can be compared with one another as are unlike, and indeed the example
of disparity adduced in §.~26. can be obtained in this way; so that all that has hitherto been handed
down is contained in this general investigation. Let there be feigned, namely, this integral
equation:
\[ \alpha xxyy + 2\beta xxy + 2\gamma xyy + \delta xx + \varepsilon yy + 2\zeta xy + 2\eta x + 2\theta y + \varkappa = 0 \tag{1} \]
from which there arises
\[ y = \frac{-\beta xx - \zeta x - \theta + \sqrt{(\beta xx + \zeta x + \theta)^{2} - (\alpha xx + 2\gamma x + \varepsilon)(\delta xx + 2\eta x + \varkappa)}}{\alpha xx + 2\gamma x + \varepsilon} \tag{2} \]
\[ x = \frac{-\gamma yy - \zeta y - \eta - \sqrt{(\gamma yy + \zeta y + \eta)^{2} - (\alpha yy + 2\beta y + \delta)(\varepsilon yy + 2\theta y + \varkappa)}}{\alpha yy + 2\beta y + \delta} \tag{3} \]
Let there now be put, for brevity's sake:
\[ \begin{aligned}
App &= \beta\beta - \alpha\delta &\qquad \mathfrak{A}qq &= \gamma\gamma - \alpha\varepsilon \\
2Bpp &= 2\beta\zeta - 2\alpha\eta - 2\gamma\delta &\qquad 2\mathfrak{B}qq &= 2\gamma\zeta - 2\alpha\theta - 2\beta\varepsilon \\
Cpp &= \zeta\zeta + 2\beta\theta - \alpha\varkappa - \delta\varepsilon - 4\gamma\eta &\qquad \mathfrak{C}qq &= \zeta\zeta + 2\gamma\eta - \alpha\varkappa - \delta\varepsilon - 4\beta\theta \\
2Dpp &= 2\zeta\theta - 2\gamma\varkappa - 2\varepsilon\eta &\qquad 2\mathfrak{D}qq &= 2\zeta\eta - 2\beta\varkappa - 2\delta\theta \\
Epp &= \theta\theta - \varepsilon\varkappa &\qquad \mathfrak{E}qq &= \eta\eta - \delta\varkappa
\end{aligned} \]
and there will be:
\[ p\sqrt{Ax^{4} + 2Bx^{3} + Cxx + 2Dx + E} = \alpha xxy + 2\gamma xy + \varepsilon y + \beta xx + \zeta x + \theta \tag{4} \]
\[ -q\sqrt{\mathfrak{A}y^{4} + 2\mathfrak{B}y^{3} + \mathfrak{C}yy + 2\mathfrak{D}y + \mathfrak{E}} = \alpha xyy + 2\beta xy + \delta x + \gamma yy + \zeta y + \eta \tag{5} \]

\origpage{55}
\textbf{§.~33.} But if the assumed integral equation be differentiated, there will result
\[ dx(\alpha xyy + 2\beta xy + \gamma yy + \delta x + \zeta y + \eta)
   + dy(\alpha xxy + \beta xx + 2\gamma xy + \varepsilon y + \zeta x + \theta) = 0 \tag{6} \]
whence, if the values (4) and (5) found for those factors be substituted, there will arise this
differential equation:
\[ \frac{q\,dx}{\sqrt{Ax^{4} + 2Bx^{3} + Cxx + 2Dx + E}}
   = \frac{p\,dy}{\sqrt{\mathfrak{A}y^{4} + 2\mathfrak{B}y^{3} + \mathfrak{C}yy + 2\mathfrak{D}y + \mathfrak{E}}} \tag{7} \]
the integral of which is therefore the assumed equation (1). But since above there are 10 equations,
while the number of the coefficients $\alpha$, $\beta$, $\gamma$, $\delta$, etc. is 9, of which one
can be assumed at pleasure, eight letters will remain to be determined. Further, there come in
addition two letters $p$ and $q$ to be defined, so that now there are ten unknown quantities present;
from which the coefficients of each formula $A$, $B$, $C$, $D$, $E$ and $\mathfrak{A}$, $\mathfrak{B}$,
$\mathfrak{C}$, $\mathfrak{D}$, $\mathfrak{E}$ seem able to be assumed at pleasure. But it is evident
that, since the one set have already been assumed at will, the others do not depend entirely on our
choice, for otherwise any formula whatever could be reduced to an algebraic one.

\textbf{§.~34.} But hence other not inelegant transmutations of a given formula can be obtained, if
in place of $y$ other values be substituted. For instance, if $\mathfrak{C} = 0$ be set,\ednote{The text sets $\mathfrak C = 0$, but the gloss $\eta\eta = \delta\varkappa$ is in fact the condition $\mathfrak E = 0$ (since $\mathfrak E qq = \eta\eta - \delta\varkappa$); equation (8) that follows holds only when $\mathfrak E = 0$. An apparent misprint for $\mathfrak E = 0$, reproduced here as printed.} that is
$\eta\eta = \delta\varkappa$, and $y = zz$ be posited, the following differential equation will come
out.
\[ \frac{q\,dx}{\sqrt{Ax^{4} + 2Bx^{3} + Cxx + 2Dx + E}}
   = \frac{2p\,dz}{\sqrt{\mathfrak{A}z^{6} + 2\mathfrak{B}z^{4} + \mathfrak{C}z^{2} + 2\mathfrak{D}}} \tag{8} \]
the integral of which is therefore the assumed equation, if $y = zz$ be posited, and
$\eta\eta = \delta\varkappa$ be set, and the remaining letters be duly
\origpage{56}
determined. The complete integral too will be found with no difficulty; for even if perchance the
integral found does not involve a new constant, let there be put
\[ \frac{q\,dx}{\sqrt{Ax^{4} + 2Bx^{3} + Cxx + 2Dx + E}} = \frac{q\,du}{\sqrt{Au^{4} + 2Bu^{3} + Cuu + 2Du + E}} \]
and the complete integral of this equation may be assigned from what precedes; and hence the
complete integral too of the equation composed of unlike formulas will be gathered.

\textbf{§.~35.} Just as of this differential equation — to begin from the simplest:
\[ \frac{dx}{\sqrt{f + gx}} = \frac{dy}{\sqrt{f + gy}} \]
the complete integral is:
\[ gg(xx + yy) - 2ggxy - 2ccg(x+y) + c^{4} - 4ccf = 0 \]
Then, of this differential equation
\[ \frac{dx}{\sqrt{f + gxx}} = \frac{dy}{\sqrt{f + gyy}} \]
the complete integral is:
\[ xx + yy - 2xy\sqrt{1 + fgcc} - ccff = 0 \]
Thirdly, of this differential equation
\[ \frac{dx}{\sqrt{f + gx^{3}}} = \frac{dy}{\sqrt{f + gy^{3}}} \]
the complete integral is
\[ f(xx + yy) + \frac{ggcc}{4f}xxyy - gcxy(x+y) - 2fxy - gcc(x+y) - 2fc = 0 \]
Fourthly, further, of this differential equation
\[ \frac{dx}{\sqrt{f + gx^{4}}} = \frac{dy}{\sqrt{f + gy^{4}}} \]
the complete integral has been found
\[ f(xx + yy) - fcc - gccxxyy - 2xy\sqrt{f(f + gc^{4})} = 0 \]

\origpage{57}
So too the complete integral of this equation
\[ \frac{dx}{\sqrt{f + gx^{6}}} = \frac{dy}{\sqrt{f + gy^{6}}} \]
will be able to be found:

\textbf{§.~36.} Let there first be determined, in §.~33., the values, so that this equation may come
out
\[ \frac{dx}{\sqrt{fx + gx^{4}}} = \frac{dy}{\sqrt{fy + gy^{4}}} \]
whose complete integral is found:
\[ gg(xx + yy) - 4ggcxxyy - 4fgccxy(x+y) - 2ggxy - 2fgc(x+y) + ffcc = 0 \]
Let there now be put $x = tt$ and $y = uu$, so that this differential equation may come out
\[ \frac{dt}{\sqrt{f + gt^{6}}} = \frac{du}{\sqrt{f + gu^{6}}} \]
the complete integral of which will therefore be
\[ gg(t^{4} + u^{4}) - 4ggct^{4}u^{4} - 4fgccttuu(tt+uu) - 2ggttuu - 2fgc(tt+uu) + ffcc = 0 \]
whence the case resulting from the hypothesis $c = \infty$ deserves to be noted, which gives
$4gttuu(tt+uu) = f$.

\end{document}
